% ===== PRÉAMBULE CANONIQUE — à coller VERBATIM en tête de CHAQUE .tex =====
\documentclass[11pt,a4paper]{article}
\usepackage[utf8]{inputenc}\usepackage[T1]{fontenc}\usepackage{lmodern}
\usepackage[french]{babel}
\usepackage{amsmath,amssymb,mathtools}\usepackage{icomma}
\usepackage{siunitx}\sisetup{locale=FR}  % \qty{}{}, \si{}, \ang{}
\usepackage{xcolor}\usepackage{colortbl}
\usepackage{tikz}\usepackage{pgfplots}\pgfplotsset{compat=1.18}
\usepackage[siunitx,europeanresistors]{circuitikz} % circuits (élec.)
\usepackage[most]{tcolorbox}\usepackage{enumitem}\usepackage{array,booktabs}
\usepackage[a4paper,margin=2.2cm]{geometry}
\usetikzlibrary{arrows.meta,calc,angles,quotes,positioning,patterns,%
  backgrounds,shapes.geometric,decorations.pathmorphing,decorations.markings,intersections}
\definecolor{primary}{RGB}{222,49,99}\definecolor{primarydark}{RGB}{150,22,55}
\definecolor{defcol}{RGB}{0,114,178}\definecolor{methcol}{RGB}{52,92,148}
\definecolor{excol}{RGB}{0,135,90}\definecolor{attncol}{RGB}{200,40,55}
\tcbset{boxbase/.style={enhanced,breakable,boxrule=0.9pt,arc=2.5pt,
  left=3mm,right=3mm,top=2.3mm,bottom=2.3mm,fonttitle=\bfseries,coltitle=white}}
% --- boîtes TUTORIEL (noms EXACTS, ne pas renommer) ---
\newtcolorbox{cle}[1][]{boxbase,colback=defcol!6,colframe=defcol,
  title={\textsc{À retenir}\ifx\relax#1\relax\relax\else\textmd{ : \itshape #1}\fi}}
\newtcolorbox{methode}[1][]{boxbase,colback=methcol!7,colframe=methcol,sharp corners,
  title={\textsc{Méthode}\ifx\relax#1\relax\relax\else\textmd{ --- \itshape #1}\fi}}
\newtcolorbox{exemplebox}[1][]{boxbase,colback=excol!5,colframe=excol,
  title={\textsc{Exemple}\ifx\relax#1\relax\relax\else\textmd{ : \itshape #1}\fi}}
\newtcolorbox{piege}[1][]{boxbase,colback=attncol!6,colframe=attncol,
  title={\textsc{Piège}\ifx\relax#1\relax\relax\else\textmd{ : \itshape #1}\fi}}
\newtcolorbox{remarque}[1][]{boxbase,colback=black!4,colframe=black!45,
  title={\textsc{Remarque}\ifx\relax#1\relax\relax\else\textmd{ : \itshape #1}\fi}}
% --- boîtes EXERCICE / CORRIGÉ (noms EXACTS, auto-comptées) ---
\newtcolorbox[auto counter]{exo}[1][]{boxbase,colback=primary!4,colframe=primary,
  title={\textsc{Exercice}~\thetcbcounter\ifx\relax#1\relax\relax\else\textmd{ --- \itshape #1}\fi}}
\newtcolorbox[auto counter]{corrige}[1][]{boxbase,colback=excol!4,colframe=excol,
  title={\textsc{Corrigé}~\thetcbcounter\ifx\relax#1\relax\relax\else\textmd{ --- \itshape #1}\fi}}
\newcommand{\vv}[1]{\vec{#1}}\newcommand{\ds}{\displaystyle}
\newcommand{\R}{\mathbb{R}}\newcommand{\N}{\mathbb{N}}\newcommand{\Z}{\mathbb{Z}}
% (un \usepackage{fancyhdr}+en-tête est possible ; il est ignoré côté web)
% =========================================================================
\begin{document}
\sisetup{per-mode=symbol}

\begin{corrige}[distance d'arrêt et carré de la vitesse]
\begin{enumerate}[label=\alph*)]
  \item À l'arrêt $v=0$, Torricelli donne $0=v_0^2+2a\,\Delta x$, soit $\Delta x=\dfrac{v_0^2}{2|a|}$ : à freinage $|a|$ fixé, $\Delta x\propto v_0^2$.
  \item $\Delta x_2=\Delta x_1\left(\dfrac{v_2}{v_1}\right)^2=25\times\left(\dfrac{70}{50}\right)^2=25\times1{,}96=\qty{49}{\metre}$.
  \item $25\times\left(\dfrac{100}{50}\right)^2=25\times4=\qty{100}{\metre}$ : doubler la vitesse quadruple la distance d'arrêt.
\end{enumerate}
\end{corrige}

\begin{corrige}[éviter l'obstacle à temps]
\begin{enumerate}[label=\alph*)]
  \item $|a|_{\min}=\dfrac{v_0^2}{2\,d}=\dfrac{80^2}{2\times1000}=\dfrac{6400}{2000}=\qty{3,2}{\metre\per\second\squared}$.
  \item $t=\dfrac{v_0}{|a|}=\dfrac{80}{3{,}2}=\qty{25}{\second}$.
  \item Non : pendant la \qty{1}{\second} de réaction, le train parcourt $80\times1=\qty{80}{\metre}$ à vitesse constante ; il ne reste que \qty{920}{\metre} pour freiner, ce qui exige $|a|=\dfrac{80^2}{2\times920}\approx\qty{3,5}{\metre\per\second\squared}>\qty{3,2}{\metre\per\second\squared}$. Le freinage serait insuffisant.
\end{enumerate}
\end{corrige}

\begin{corrige}[deux véhicules se croisent]
\begin{enumerate}[label=\alph*)]
  \item Ils se rapprochent à $90+70=\qty{160}{\kilo\metre\per\hour}$ : $t=\dfrac{80}{160}=\qty{0,5}{\hour}$ (30 min).
  \item Distance depuis $A$ : $90\times0{,}5=\qty{45}{\kilo\metre}$.
  \item Second véhicule : $70\times0{,}5=\qty{35}{\kilo\metre}$ depuis $B$ ; $45+35=\qty{80}{\kilo\metre}$ $\checkmark$.
\end{enumerate}
\end{corrige}

\begin{corrige}[la poursuite du policier]
$v_{\text{auto}}=\qty{72}{\kilo\metre\per\hour}=\qty{20}{\metre\per\second}$.
\begin{enumerate}[label=\alph*)]
  \item Automobiliste (MRU) : $x_1(t)=20\,t$. Policier (MRUA, repos) : $x_2(t)=\tfrac12\times5\times t^2=2{,}5\,t^2$.
  \item Rencontre : $2{,}5\,t^2=20\,t \Rightarrow t=\dfrac{20}{2{,}5}=\qty{8}{\second}$ ($t\neq0$).
  \item $x_2(8)=2{,}5\times64=\qty{160}{\metre}$ ; vitesse $v=at=5\times8=\qty{40}{\metre\per\second}$ ($=\qty{144}{\kilo\metre\per\hour}$, soit le double de la vitesse de l'automobiliste).
\end{enumerate}
\end{corrige}

\begin{corrige}[la Formule 1 et le doublement du temps]
$v_0=\qty{216}{\kilo\metre\per\hour}=\qty{60}{\metre\per\second}$.
\begin{enumerate}[label=\alph*)]
  \item $v=v_0+at=60-5\times6=\qty{30}{\metre\per\second}>0$ : elle roule encore, le freinage dure bien \qty{6}{\second}.
  \item $\Delta x=v_0t+\tfrac12 at^2=60\times6-\tfrac12\times5\times36=360-90=\qty{270}{\metre}$.
  \item $x=\tfrac12 a t^2\propto t^2$ : en doublant le temps, la distance est multipliée par $2^2=\mathbf{4}$.
\end{enumerate}
\end{corrige}

\end{document}
